
\chapter{The snap, the rate, and the one asymptotic advantage}

\textbf{Purpose:} Build the fourth term as the operation it exists for, measure what a twist and a
relaxation actually move, and locate the single place in this work where a quantum algorithm holds a
genuine asymptotic advantage. \textbf{Scope:} \texttt{snap\_release.py},
\texttt{graph\_machine.py}, \texttt{hhl\_nested.py}, \texttt{product\_machine.py}.

\section{The wrong question, asked first}

Static capacity is how many reals it takes to write an arbitrary pure state, which is
\(2^{n+1} - 2\), against a fixed budget of \(cn\):

\begin{center}
\begin{tabular}{rrrl}
\toprule
qubits & reals needed & 4 reals gives & \\
\midrule
1 & 2 & 4 & exact \\
2 & 6 & 8 & exact \\
3 & 14 & 12 & short \\
16 & 131\,070 & 64 & short \\
\bottomrule
\end{tabular}
\end{center}

Four reals a qubit is exact through two qubits and short from three on, and raising 64 bits to 128
does not move that line by one qubit. It is a counting result, and it is not what the fourth term is
for.

The fourth term is not storage for an amplitude. It is a magnitude of 0 or 1 saying whether the qubit
is bound to anything, and its purpose is to make the snap available. \textbf{A three-term qubit
cannot ask the question the snap answers}: with a direction and a polarization, a qubit carries where
it points, and nothing in it says whether it is bound. The fourth term is the predicate the operation
is conditioned on.

\section{The two snaps}

\textbf{Release}, a \(Z\) measurement: delete every edge on the qubit. It leaves the cluster and its
entanglement is gone. \textbf{Switch}, a \(Y\) measurement: complement the edges \emph{among} the
qubit's neighbors, then delete its own edges. The neighbors inherit the connectivity, and the
entanglement moves, and does not vanish.

Order is not interchangeable. The complementation must happen while the qubit still has its edges,
because the edges define the neighborhood. A delete first would complement an empty set, and switch
would silently degrade into release.

\section{The exact null on the twist}

A twist is local complementation, and local complementation is a local Clifford operation. Those
preserve entanglement across every bipartition. So:

\begin{quote}
The entanglement response to a twist is \textbf{exactly \(0.0\)} at every vertex, on every shape, at
any size.
\end{quote}

Computed, and not remembered. That is an exact null in the same family as the others here, and it
says the thing this tree keeps finding from new directions: \textbf{relabelling is free.} A twist
moves the description and not the content.

The edge count is not a local Clifford invariant. The tension therefore moves while the correlation does
not. That is the torsion against tension relationship, with the conservation measured, and not
asserted, and the tension moment separates all four test shapes at \(0.75\), \(1.0\), \(55.125\) and
\(441.0\).

\section{Dissipation, which reads null}

The model has no dissipation in it at all: local complementation and measurement are closed, nothing
leaks and nothing relaxes. Dissipation therefore cannot explain the tension table, and the table is
neighborhood parity.

Added as a process: release qubits one at a time until nothing is bound. \textbf{Driven} takes the
qubit under the most neighbor pressure, recomputed each step, where pressure is the sum of the
\emph{neighbors'} degrees, and not the qubit's own. \textbf{Random} takes a fixed order over 20
seeds and is the matched control.

\begin{center}
\begin{tabular}{lrrrl}
\toprule
shape & driven rate & random mean & random range, 20 seeds & verdict \\
\midrule
line & 0.3612 & 0.3248 & 0.2785 -- 0.4494 & inside \\
ring & 0.3006 & 0.3251 & 0.2473 -- 0.4056 & inside \\
complete & 0.0501 & 0.0501 & 0.0501 -- 0.0501 & inside \\
grid 2x5 & 0.3583 & 0.3006 & 0.2397 -- 0.4094 & inside \\
fibonacci & 0.2617 & 0.2573 & 0.2272 -- 0.2924 & inside \\
golden sphere & 0.2747 & 0.2692 & 0.2310 -- 0.3140 & inside \\
\bottomrule
\end{tabular}
\end{center}

\textbf{Zero of six shapes put the driven rate outside the random orders' own range.} Neighbor
pressure changes nothing measurable: the relaxation rate is a property of the shape and not of the
order of release. The complete-graph row is a free exact control, since every qubit in it carries
identical neighbor pressure and the range has width exactly zero.

What remains is the rate itself, intrinsic to the topology and robust to order. A complete graph
resists relaxation about \(6\times\) harder than a line, with half-lives of 13.85 steps against 1.92.

\section{A gap needs the range of its control}

Gaps of \(+0.0457\), \(-0.0200\) and \(+0.0641\) against rates near \(0.3\), counted as ``differed on 5
of 6'' by exact floating-point inequality, would call pressure a real driver. They have no null: without
the spread from seed to seed of the random control, no gap can show that it exceeds noise. Gaps of both
signs alone forbid the claim. Against the control's own range over 20 seeds, every gap sits inside it.

\section{The one genuine asymptotic quantum win}

HHL solves \(Ax = b\) in \(O(\log(N)\, s^2 \kappa^2 / \epsilon)\). It needs sparsity, size, and
conditioning. A 1D qubit chain is tridiagonal, and \(s = 2\), and the graph representation reaches a
million qubits in 8 MB. Conditioning is the only open question and the diagonal decides it.

A connected graph's Laplacian is \textbf{singular}: the all-ones vector is in its kernel. Add a
diagonal field term \(m\) and the spectrum moves from \([0, 4d]\) to \([m, 4d+m]\). Then
\(\kappa = 1 + 4d/m\), \textbf{constant in \(N\)}. Measured against the exact finite-width prediction
to \(7.28 \times 10^{-16}\), climbing monotonically to its bound through \(4.8478\), \(4.9616\),
\(4.9904\), \(4.9976\), \(4.9994\), \(4.9998\).

The field term is the lever, and the diagonal is the fourth term: a per-qubit entanglement magnitude
\emph{is} a diagonal contribution. The four-term specification and HHL's applicability are therefore the
same fact.

\begin{center}
\begin{tabular}{rrr}
\toprule
field \(m\) & \(\kappa\) & qubits needed \\
\midrule
0.25 & 17.00 & 16.00 MB \\
1.00 & 5.00 & 2.00 MB \\
8.00 & 1.50 & \textbf{0.25 MB} \\
\bottomrule
\end{tabular}
\end{center}

A \(32\times\) stronger field cuts the qubit requirement \(64\times\). More carried per qubit, fewer
qubits, and that is the conditioning bound, and not a metaphor.

Uneven field distance is punished by its weakest qubit, since
\(\kappa \le (\lambda_{\max} + m_{\max})/m_{\min}\):

\begin{center}
\begin{tabular}{lrrrr}
\toprule
shape & \(m_{\min}\) & \(m_{\max}\) & \(\kappa\) & crossover \\
\midrule
uniform & 1.000 & 1.000 & 5.00 & \(6.554 \times 10^{4}\) \\
spikeball & 0.250 & 8.000 & \textbf{40.78} & \(2.097 \times 10^{6}\) \\
fuzzball & 0.754 & 1.271 & \textbf{5.28} & \(6.554 \times 10^{4}\) \\
\bottomrule
\end{tabular}
\end{center}

A fuzzball costs almost nothing against uniform. A spikeball is \(32\times\) worse despite
carrying qubits at eight times the field.

Nesting does not lower the bar. Against the honest classical baseline, conjugate gradient at
\(O(N s \sqrt{\kappa})\), and not a dense solve, the crossover \emph{rises} with dimension from
\(6.554 \times 10^{4}\) qubits in 1D to \(8.389 \times 10^{6}\) in 8D, because \(s\) and \(\kappa\)
both grow with \(d\) and enter the quantum cost squared while the classical cost takes them linearly
and as a square root. What nesting preserves is the logarithmic \(N\) term at every dimension.

\textbf{This is a real exponential-in-\(N\) advantage, and it is on the qubit map, and not on the hash.}
It solves a linear system. SHA-256 is not one.

\section{The null permutation on placement}

The bound \(\kappa \le (\lambda_{\max} + m_{\max})/m_{\min}\) depends only on the \emph{multiset} of
field distances. A permutation of which qubit carries which value is therefore a null permutation: the object
changes and the bound cannot move.

Across 44 arrangements of one multiset the bound took \textbf{exactly one value}, \(48.0000\), while
the measured \(\kappa\) ranged from \(40.5676\) to \(47.5125\). So the bound is blind to something
real, and placement carries a factor of \(1.170\) in conditioning for the same field strength, free.

The magnitude is honest about itself. That factor is 17\% and not \(32\times\), and it does not rescue
a spikeball; even placement recovers a sixth of one factor of the thirty-two that uniformity is
worth. And the best random draw at \(40.5676\) edged the evenly spread arrangement at \(40.6150\). The
placement optimum is therefore a real problem, and not a formula.
